// Custom Blend — overlay blend mode using user-defined functions
//
// Demonstrates defining reusable functions for image compositing.
// Implements the Photoshop Overlay blend mode per-channel, then
// mixes the result with the original base using $blend_amount.
//
// NOTE: TEX now ships built-in overlay()/soft_light()/screen()/... that work
// directly on vec3 (see the Blend section of the cookbook). This example keeps
// its own my_overlay()/my_soft_light() to show how to write a CUSTOM blend —
// user functions may NOT reuse a built-in's name, so we prefix them with my_.
//
// Each function tests a pixel value, so its `if` is per-pixel: both branches
// run on every pixel and the result is merged pixel by pixel. A `return`
// inside such an `if` would return for EVERY pixel (LANGUAGE.md §7.1), so each
// function assigns its result inside the `if` and returns once at the end.
//
// Connect @base and @overlay to RGB images.

float my_overlay(float a, float b) {
    float res = 1.0 - 2.0 * (1.0 - a) * (1.0 - b);
    if (a < 0.5) {
        res = 2.0 * a * b;
    }
    return res;
}

float my_soft_light(float a, float b) {
    float d = a;
    if (a < 0.25) { d = ((16.0 * a - 12.0) * a + 4.0) * a; }
    else { d = sqrt(a); }
    float res = a + (2.0 * b - 1.0) * (d - a);
    if (b < 0.5) {
        res = a - (1.0 - 2.0 * b) * a * (1.0 - a);
    }
    return res;
}

f$blend_amount = 0.5;   // mix strength (0 = base only, 1 = full overlay blend)

vec3 base = @base.rgb;
vec3 over = @overlay.rgb;

vec3 result = vec3(
    my_overlay(base.r, over.r),
    my_overlay(base.g, over.g),
    my_overlay(base.b, over.b)
);

// The built-in would be the one-liner:  vec3 result = overlay(base, over);

@OUT = lerp(base, result, $blend_amount);
